\section{Why Shortcuts Emerge}
This section explains why shortcuts emerge in complex systems and sets up the control triangle in Section 6, while foreshadowing the axes and legitimacy analysis in Sections 12 and 13. Continuous debate becomes infeasible or undesired at scale, and disputes are often bypassed rather than resolved. Luhmann frames this as a problem of complexity reduction within social systems, while Hirschman shows how exit and voice redirect conflict away from direct contestation; Lukianoff and Schlott describe contemporary cancellation dynamics that similarly short‑circuit debate rather than carry it to resolution.\autocite{luhmann1995,hirschman1970,lukianoffschlott2023}

The drivers begin with cognitive and informational constraints. Epistemic load names the complexity of contested knowledge as more issues become simultaneously debatable; Simon describes how information‑rich environments intensify these cognitive limits, and cognitive load and overload research clarifies why those limits bind under scale.\autocite{simon1971,sweller1988,epplermengis2004} These cognitive pressures are complemented by deliberative and negotiation burdens that make sustained dispute resolution costly in practice.

Deliberative load captures the time and effort cost of sustained argument even when disputes are tractable in principle; decision‑theoretic and collective‑action accounts show how bounded rationality, informational costs, and coordination costs constrain sustained participation.\autocite{simon1947,downs1957,olson1965} Negotiation load extends this logic to the social and political cost of re‑opening settlements, where transaction‑cost and collective‑choice frameworks emphasize how renegotiation frictions and decision rules make re‑settlement expensive and slow.\autocite{coase1937,williamson1985,williamson2000,buchanantullock1962}

These loads are amplified by resource constraints and issue load because attention and bandwidth limits concentrate decision‑making in saturated environments; Simon’s attention‑scarcity argument and subsequent overload studies make this concentration explicit, and Zuboff shows how it intensifies in digital environments.\autocite{simon1971,epplermengis2004,zuboff2018} External shocks compress time horizons and raise perceived stakes, accelerating the turn to shortcuts; work on states of exception and COVID‑era science politics illustrates this compression in practice.\autocite{agamben2005,jasanoff2021}

Legibility further constrains deliberation because what is illegible is more easily bypassed or excluded. Scott shows how legibility schemes filter what can be administered and thus what can be deliberated; classification and infrastructure studies show how embedded systems make some realities easier to ignore.\autocite{scott1998,bowkerstar1999,star1999} Foucault’s analysis of disciplinary visibility shows how legibility is produced and used as a control mechanism.\autocite{foucault1977} When an issue is treated as settled—self‑evident, orthodox, or dogmatic—willingness to bypass deliberation rises, and adoption becomes decisive; this bears directly on later claims about legitimacy and justification.\autocite{foucault1977}

These pressures motivate shortcut behavior. Shortcuts are intended stability‑seeking mechanisms under complexity; Luhmann emphasizes stabilization rather than truth as the system‑level function of such mechanisms, while Simon shows how bounded rationality favors satisficing shortcuts over continuous renegotiation.\autocite{luhmann1995,simon1947} As a result, shortcuts stabilize epistemic preferences by narrowing the space for renegotiation, a move Luhmann’s account of stabilized meaning helps explain.\autocite{luhmann1995} This link between cognitive limits and stabilization sets up the control model that follows.

Different groups stabilize different preferred states, and clashes emerge when these stabilizations diverge. Mouffe’s agonistic model treats such conflict as constitutive rather than resolvable by consensus, and Iyengar and Westwood provide empirical evidence of polarization that aligns with these divergent stabilizations.\autocite{mouffe2000,iyengarwestwood2014}

Power shapes the adoption and operation of shortcuts. Lukes’ account of power’s dimensions shows how power can bypass debate or marginalize dispute rather than resolve it, and Foucault’s analysis of disciplinary mechanisms shows how power can act directly on conduct without relying on persuasion.\autocite{lukes2005,foucault1977} Power also varies by form— institutional, economic, technical, and cultural— and its distribution shapes perceived legitimacy, a point returned to in the axes and legitimacy sections.\autocite{lukes2005} Platform and infrastructural power further enable low‑visibility control at scale, as analyses of opaque systems and embedded classifications indicate.\autocite{pasquale2015,bowkerstar1999,star1999}

The control triangle is a model that isolates two recurring pathways—constraint and attraction—acting on conduct in epistemic control. Wiener’s feedback framework provides the conceptual basis for this control model.\autocite{wiener1948} It ties the preceding account of shortcut pressures to the mechanisms that follow, and its scope applies to cases where epistemic preferences are stabilized through expression and visibility and through internalized norms; this covers the taxonomy’s operational modes developed later. The model is an abstraction rather than an exhaustive account, and other forms of power can act directly on conduct or bypass the triangle.
